% ============================================================
% CAPÍTULO 6
% ============================================================
\capitulo{Conclusiones}

\section{Cumplimiento de los Objetivos}
\subsection{Conjunto de Datos}
Se conformó un conjunto propio de 274 videos y 393 clips de equinos, de los cuales el sistema utiliza 359, correspondientes a 133 caballos, con un criterio de etiquetado documentado y referido a EquiFACS.

\subsection{Preprocesamiento}
Se implementó un \textit{pipeline} reproducible que estima la pose, elige una ventana temporal medida, selecciona 16 fotogramas, alinea la cara y genera representaciones visuales y geométricas, con un aprovechamiento del 91,3~\% de los clips.

\subsection{Clasificación}
El modelo MIL multimodal alcanza un \fmacro{} de 0,740 y una exactitud de 0,744 con evaluación agrupada por caballo, por encima del criterio de aceptación fijado de antemano.

\subsection{Aplicación}
Se implementó una aplicación de inferencia que usa el mismo código que el entrenamiento, verificado por pruebas automatizadas, y que se niega a clasificar cuando el video no permite una medición suficiente.

\section{Conclusión sobre la Hipótesis General}
La hipótesis general se considera respaldada dentro de los límites del conjunto: es factible clasificar automáticamente las tres expresiones definidas a partir de videos cortos, con generalización a caballos no observados y un rendimiento por encima del criterio establecido. La afirmación se refiere a las categorías y al criterio de una etiquetadora; no constituye una validación clínica ni etológica externa.

\section{Conclusión sobre la Hipótesis Temporal}
La hipótesis temporal no fue demostrada. Los modelos sensibles al orden no superaron a los invariantes, y un resumen estático de los 16 fotogramas alcanza el mismo rendimiento que el modelo final. Utilizar varios fotogramas sí resultó útil; ordenarlos o ponderarlos, no.

\section{Aportes}
\begin{itemize}
  \item Un conjunto de datos propio de expresiones faciales equinas con criterio de etiquetado documentado.
  \item La identificación y el control de la fuga de información por identidad.
  \item Un método de selección de fotogramas que mide antes de elegir y más que duplicó el aprovechamiento de los clips.
  \item Un sistema de líneas de base que permite atribuir el rendimiento a sus fuentes y detectar atajos.
  \item La equivalencia verificada entre el procesamiento de entrenamiento y el de inferencia.
  \item El registro de resultados negativos y de variantes descartadas.
\end{itemize}

\section{Trabajo Futuro}
\begin{enumerate}
  \item Medir el acuerdo entre observadores sobre una muestra del conjunto: que una segunda persona etiquete a ciegas unos 60 clips con el criterio documentado y se calcule el $\kappa$ de Cohen \parencite{cohen1960}.
  \item Incorporar más caballos, en particular de \clase{Disgustado} en contextos distintos de la manada.
  \item Validar el modelo con videos de otras fuentes y establecimientos.
  \item Equilibrar el contexto de filmación entre clases o controlarlo explícitamente.
  \item Estudiar representaciones temporales sólo si nuevos datos muestran señal temporal que el esquema actual no captura.
  \item Reducir el costo de la estimación de pose en inferencia.
\end{enumerate}

\section{Conclusión Final}
Un sistema de clasificación de expresiones faciales equinas puede construirse con un \textit{pipeline} reproducible y evaluarse con rigor evitando la fuga por identidad. El mejor resultado no provino de aumentar la complejidad del modelo sino de mejorar el aprovechamiento de los datos, la coherencia metodológica y el control de sesgos. La información de varios fotogramas resulta útil, pero con los datos y escalas estudiados no se encontró evidencia de que su orden sea determinante. El trabajo futuro debería orientarse a ampliar y diversificar el conjunto y a medir la confiabilidad de sus etiquetas antes que a arquitecturas temporales más complejas.
